% Generated by records_reorg.py assemble (records-reorg/); edit the data and re-run, not this file.
\subsection{Principal limitations}\label{recsec:limitations-principal}

% Part 6 chunk 3 (the director's compression at depth, 2026-11-20; DECISION-20261120a item 5), director-helper-3: each result's limitations, moved verbatim from sec_limitations.tex @ 940561ea6 lineage (H3-P6-10..15). References into the note resolve through main_note.aux (xr). The note's per-result pointer subsections were dropped in the pointer pass (journal-editor, 2026-11-20, JE-24); this section's subsections carry their own labels (rlim:...). [records reorg] former heading: \section{Limitations moved from the note (Part 6)} [records reorg, not printed] This section holds each result's limitations as the note stated them before Part 6 was compressed. The note keeps each result's headline limits in its row of Table~\ref{tab:validity}; its \S\ref{sec:limitations} points to this section, and the body pointers that once led to its per-result limitation subsections now name the result's subsection here.

\paragraph{Principal limitations.} The body states each limitation with the result it qualifies; the principal ones are these. $C_k$ denotes coverage within $\pm k$ standard deviations. The quoted 2D statistical block has had no coverage test. The fixed-truth coverage studies scored two earlier statistical bands, of which the first-built undercovers in its tails ($C_2=\covTwoDCTwo$, vs.\ $0.954$; \S\ref{sec:uq-boot}, \S\ref{sec:coverage}). No result's intervals have validated coverage. In two dimensions, the procedure's covariance under-covers jointly at both held-out tunes, with \lrval{cov2d-ar23-joint} of \lrval{cov2d-ntoys} pseudo-experiments in the joint region at each (\S\ref{sec:coverage-tunes}, \S\ref{rlim:2d}). The five-dimensional coverage test is in \S\ref{sec:coverage-5d}, and the sixteen cells' second, descriptive coverage test in \S\ref{sec:fived16-fd} and Table~\ref{tab:test-index}. The other principal limitations are stated with their results:
% [records reorg] former heading: \subsection{Principal limitations (from \S\ref{sec:limitations})}
\begin{itemize}
\item The generator comparisons (\S\S\ref{sec:3d}--\ref{sec:eavailw}).
\item The absence of a 3D or 4D covariance (\S\ref{sec:3d}).
\item The five-dimensional estimator's regularization bias (\S\ref{sec:val-departure}) and the purity-subtraction bias of the \estPurity{} and the \estDetPurity{} (\S\ref{sec:val-negweight}).
\item Five-dimensional intervals and their coverage (\S\ref{sec:gbdt-model-dependence}, \S\ref{sec:coverage-5d}).
\item The five-axis covariance of the \resFiveAxis{} (\S\ref{sec:eavailw-adopted}).
\item The calibration and nuisance model of the generator tests (\S\ref{sec:joint5d-frozen}, the earlier note's section on the hadronic-response condition).
\item Tier~2 of the extended-fiducial study (App.~\ref{sec:fps}).
\item The fake-data limitation of the sixteen-cell result (\S\ref{sec:fived16-fd}).
\item The dependence of the sixteen cells on the unfolding algorithm (\S\ref{sec:fived16-ibu}).
\item The prior retained at the two truths flagged under the baseline covariance (\S\ref{sec:fived16-prr}).
\item The bias-half screen (\S\ref{sec:fived16-fd}).
\item The ensemble-averaged variant, not adopted (\S\ref{sec:val-departure}).
\item The run-period inconsistency present before unfolding (\S\ref{sec:uq-combine}, \S\ref{sec:fived16-pl}).
\item Point-cloud work (the earlier note's appendix pointing to the companion PET note).
\end{itemize}
% Register round n1 (S1; caveat-study R-N1; the model is paper_body.tex sec:limitations): the former "What is not claimed" list. Each fact of each former item prints at its home in the body; ledger-nA.md rows A* map each one. Sources of the former items: 2D coverage, OUTCOME-20261005-2d-fixed-truth-coverage-fail.md (VL169) and KNOWN_ISSUES 85 (the quoted block is S1, stageC/production/DESIGN.md sec. 3.4); 5D coverage, OUTCOME-20260925-s5c-tier-s-futility-fail.md (VL150), KNOWN_ISSUES 98, M2-RATEFAIL-DIAGNOSIS.md (CITABLE FOR the decomposition with its inferred status, NOT CITABLE FOR the source of the W-graded term); regularization bias, OUTCOME-20260925-s5n-stage1-development-fail.md (VL151-VL152), KNOWN_ISSUES 75, 77; estimator stability, KNOWN_ISSUES 80, RECORD-20260927-s5p-stage2-exit.md sec. 2 (VL155), OUTCOME-20260926-s5e-oi192-diagnosis-and-candidate.md, DIAGNOSIS-20260926-s5e-oi192-estimator.md (VL153-VL154); five-dimensional measurement, RECORD-20260927-s5p-stage2-exit.md sec. 4; background treatment, Table tab:estimators-matrix and app_negweight.tex sec:negweight-footing. Each absence claim below is restated from its home: "has had no coverage test" (sec_systematics.tex sec:uq-boot); "no 3D or 4D covariance is quoted" (sec_3d.tex lead); "enter no quoted uncertainty" (sec_5d.tex sec:val-departure); "enters no quoted covariance" (sec_5d.tex sec:val-negweight); "no five-dimensional measurement with total uncertainties was constructed" (sec:gbdt-model-dependence); "No frequentist coverage is therefore claimed" (sec:coverage-5d); "no calorimetric response band" (sec:joint5d-response).
% Search: held-out simulation; held-out-MC re-test; S1 has not been coverage-tested; coverage statement about Q_R (the search of its home, sec_systematics.tex sec:uq-boot) @ ddf16f10 Review fix (outline: the \covTwoDCTwo uses are not touched; check_dead_containment.py counts four prints). The sentence that prints it, the second of this paragraph, is unchanged in n2. Note consolidation (DECISION-20261027 item 1, cost-list R01; director's ruling: a summary points to the body home, never to another summary): the other limitations are no longer restated here; each is named once with a pointer to its body home. The 2D coverage sentence is kept unchanged (a \covTwoDCTwo use, out of scope). Of the absence claims listed above, "has had no coverage test" still prints in this paragraph, and the 3D and 4D one prints as "the absence of a 3D or 4D covariance" (stage-2 review fix: a bare topic there read as if 3D and 4D uncertainties existed).
% - Note round 1 (2026-11-08, draft): COV2D sentence, OUTCOME-20261108-cov2d-both-tunes-under-cover.md @ b551b05f :99-104 (CITABLE FOR) and :110, :115 (NOT CITABLE FOR); the frozen KI-85 sentence before it is unchanged. U16 (currency check of 2fe070fec): the findings of U1 (FD-PRR readout 9238f22b), U2 (V-HALF readout e7423390), U3 (ENS1, OUTCOME-20261109-ens1-stab-pass-bias-worse.md) and U15 (runperiod_values.json f63074b1), as the paper's abstract and introduction carry them (DECISION-20261116-joseph-paper-batch.md item 39).
The tests that limit a reported result are indexed in Table~\ref{tab:test-index}, with the criterion each was graded against.
% Moved to Part 5 (part5_validation.tex) by cov2d's landing bd79294b2: move-manifest-je-test-index.tsv JE-06.01..61, received as CV-R06.01..61 (2026-11-19).
% Part 6 (the director's Part 6 design, 2026-11-19; director-helper-3, journal-editor's ack): tab:result-status and its
% lead sentence are replaced by tab:validity (sec_validity_table.tex), which carries this label as a second label; each
% cell's home is listed in move-manifest-h3 rows H3-P6-*. The table's Search line moved with the statement it searches
% (tab:validity, the 2D row: the statistical block has had no coverage test).
% Move pair H3-014 (director-helper-3; move-manifest-h3-20261117.tsv sends): the 2D result's validation status, four
% sentences of sec_systematics.tex's "Status of the 2D result" paragraph at 2fe070fec, verbatim (COMBINED-PLAN sec. 4
% item 7), with their Search line; sec:uq-2d-scope keeps the construction and points here.
% Pointer entry 5 (the director's ruling (a), via director-helper-3; limitations-export row n05): sec_experiment.tex:40, kept there.

\paragraph{Limitations.}
% Check stage of the journal-voice rewrite (2026-11-21; CHECK-STAGE.md item 2; DECISION-20261120a item 5), director-helper-3:
% register rows whose limitation the rewritten Secs. 1, 3 and 6 no longer state. Each sentence below is moved verbatim
% from the note as it stood at e89694209 (draft/note-round2-20261108); the source file and line are given with each.
% Under DECISION-20261121b a records-status register row needs only its verbatim text here; the note carries no
% pointer to this section (Joseph, 2026-11-21: note-to-records pointers dropped).
% [records reorg] former heading: \section{Limitations moved from the note in the journal-voice rewrite}

\subsection{Limitations by result}\label{recsec:limitations-results}

\paragraph{Data and event loop.}
% [records reorg] former heading: \subsection{Data and event loop (from \S\ref{sec:limitations})}
\label{rlim:experiment}
The flux histogram is POT-weighted across playlists; playlist-to-playlist flux differences are a $\sim\SI{0.2}{\percent}$ effect (\S\ref{sec:evloop}).
% Part 6 (the director's ruling (A), 2026-11-19): the 2D block returns here as the file sec_q1_limitations.tex, moved whole
% and byte-identical from Part 3 (H3-P6-04); journal-editor's JE-07 rows stay as the record of its first landing.

\paragraph{The two-dimensional cross section.}
% [records reorg] former heading: \subsection{The two-dimensional cross section (from \S\ref{sec:limitations})}
\label{rlim:2d}
% Part 3 (Q1), landed verbatim from sec_limitations.tex:134-249 @ 3f88c127c (journal-editor's manifest JE-07.001..102,
% ~/codex-reports/note-measurement/part3/move-manifest-je-limitations-2d.tsv; the 2026-11-19 agreement): section 12's 2D block
% retires into Part 3, where its paragraphs go to the 2D uncertainty, comparison and validation sections as Part 3 is condensed.
\subsubsection{The two-dimensional cross section}
\label{sec:limitations-2d}

The result's intervals, bias and covariance have not been independently validated. In particular, the coverage of its intervals has not been validated on independent simulated samples, nor simultaneously over all reported bins and projections. At two held-out GENIE~v3 tunes, a test of the procedure's 2D covariance remade at each tune's truth, with thresholds fixed before it ran, reads under-covers, with a joint in-region count of \lrval{cov2d-ar23-joint} of \lrval{cov2d-ntoys} at both (\S\ref{sec:coverage-tunes}); it concerns the procedure's covariance at those tunes, not the reported intervals at the data. Such validation would need an independent simulation sample of production size, a validated generator of pseudo-data with continuously varied nuisance parameters, and samples from independent event generators beyond the truth-level ones already examined; none is available.
% Search: independent validation; held-out re-test; KI-85; S1 coverage; simultaneous coverage; fixed-truth test VL169; independent production-size MC bank; continuous nuisance generator; independent-generator reco sample; Stage D inputs (S1 not coverage-tested, re-test deferred, KNOWN_ISSUES.md:75; VL169 is a same-simulation test of an earlier band, sec_validation.tex:135-160; all three inputs UNAVAILABLE or unvalidated, state/2d-closeout/stageD-inputs/REPORT.md:59, :66-68; 0 disjoint production-size sets, DESIGN-20261008-2d-independent-statistical-validation.md:46) @ 0d919b1a

% Move pairs CV-B01..CV-B34 (cov2d; history-cuts/move-manifest-cv-20261118.tsv sends; journal-editor receives): the 2D
% limitations that sec_systematics.tex @79a934a41 stated and this section did not, verbatim with their source comment
% lines, except: CV-B03 adds the pointer (\S\ref{sec:uq-boot}); CV-B05 drops its lead clause "Its validation status is
% stated in \S\ref{sec:limitations}:" (a self-pointer here) and capitalizes "We"; in the third list item "(below)" becomes
% (\S\ref{sec:uq-2d-scope}). The paragraph headings are new, except "Seed-related terms reported, not added.".
% Move pairs CV-VB01..CV-VB29 (cov2d; history-cuts/move-manifest-cv-20261118.tsv sends; journal-editor receives): the
% validation limitations of sec_validation.tex @79a934a41, verbatim with their comment lines. The paragraph headings are
% new. Proposed order: the closure paragraph and the requirements paragraph after the existing first paragraph; the
% fixed-truth paragraph next to chunk 1's CV-B02/CV-B03 (CV-B03's "also" then has its antecedent); the held-out-tune
% paragraph last. CV-VB (P21 sentence 3) replaces chunk 1's CV-B01, now a dedup row.
% CV-VB01..CV-VB04 (sec_validation.tex:39; sentences 1-2 also stay in the body as the scope label):
\paragraph{The closure tests (\S\ref{sec:validation}).} These closure tests recover injected features and are not a formal statistical validation of the estimator. No non-closure uncertainty is assigned from them. Assigning one would require a suite of injection shapes and amplitudes specified before the tests are run. Calibration to a single Gaussian bump would leave that requirement unmet.

% CV-VB05..CV-VB11 (sec_validation.tex:48 with its comments :49-56; its last two sentences restate in other words the
% last sentence of this subsection's first paragraph, so journal-editor may merge the two):
\paragraph{Requirements for validated total uncertainties.} The follow-up assessment~\cite{TwoDFollowup2026} distinguishes a proposed split-sample width check from repeated-sampling coverage. With the MC sample, binning maps and background subtraction held fixed, such a check could establish only the precision conditional on them, at its tested exposure. Nonlinear unfolding does not guarantee a simple scaling from half to full exposure. Nor can the check supply the variation across independent MC populations. With the current inputs, validated total uncertainties for all 205 bins, the integral and the projections therefore cannot be obtained. They would require independent populations or a validated generative law, a continuous nuisance generator for tests of total intervals, and an untouched generator-truth domain for bias checks (\S\ref{sec:uq-2d-scope}). This is a limit of the available inputs, not a general impossibility result for 2D unfolding.
% Accuracy fix (unscored; director ruling 2026-10-26): the 2D result now uses refined negative-weight subtraction
% (sec:method-algo), so the condition names the background subtraction, not the purity correction.
% J3/J4 provisional (DECISION-20261018, packet C3; review m1): "and the construction repairs in sec:uq-2d-construction"
% removed; that section recorded that the quoted construction carries those fixes. Since note register round n1
% (batch n1-B), the fixes are stated in sec:uq-2d-scope, paragraph "Changes from the earlier 2D budget".
% Note register round n2, batch n2-B (2026-10-26; R-N2a): the first sentence after the condition is split at its
% scaling clause; no claim changes.
% Search: split-sample width check; independent MC populations; validated generative law @ 1a015eef

% CV-VB12..CV-VB17 (sec_validation.tex:204 with its comment :205; sentence 1 also stays in the body as the scope label):
\paragraph{The fixed-truth coverage test (\S\ref{sec:coverage}).} The test concerns the statistical band alone and changes no central value or quoted uncertainty. It tests the LightGBM statistical estimator at the band's training seed (seed~1), while the quoted central value uses seed~42. Systematic, total and five-dimensional coverage are outside this test's scope. Its simulation-side scatter reuses the resampling that built the band. Only the data-side fluctuation is tested by an independent mechanism. Because $T$ equals the prior, the test cannot detect regularization bias toward the prior.
% J3/J4 provisional (DECISION-20261018, packet W3: D2): true of the test; the quoted statistical block is at seed 42.
% CV-B02, CV-B03 (sec_systematics.tex:54, sentences 3, 4; CV-B01, sentence 6, is a dedup of CV-VB (P21 sentence 3)):
No coverage test on held-out simulation has been run, and no design for one has been registered in advance. Because the fixed-truth test's truth equals the prior, it also cannot decide whether the smaller statistical block (\S\ref{sec:uq-boot}) reflects an efficiency gain, regularization bias or undercoverage.
% Search: held-out simulation; held-out-MC re-test; S1 has not been coverage-tested; coverage statement about Q_R (the search of its home, sec_systematics.tex sec:uq-boot) @ ddf16f10

% CV-VB18..CV-VB29 (sec_validation.tex:234 with its comments :235-238, and :240, sentences 1-3, with a copy of its
% comments :241-250; sentence 4 stays in the body):
\paragraph{The held-out-tune coverage test (\S\ref{sec:coverage-tunes}).} The test resolves a per-cell deficit in content of about \lrval{cov2d-res-cell}, and about \lrval{cov2d-res-joint} jointly; smaller deficits can pass. It is one-sided: over-coverage is reported, not failed. The pseudo-experiments depart from the null's law in known ways: the flux and pair-band draws are not centered, clipping truncates the shifts by setting a negative composite weight to zero, separately on the reconstructed side and on the truth side, and three terms are in the remade covariance and are not drawn in the pseudo-experiments (the training-seed block, \S\ref{sec:uq-ml}, $C_{\mathrm{model}}$ and the simulation-statistics part of the statistical term). Their effect on the flag rate was not measured. The false-flag rates are those of a null in which the remade covariance is exact. That covariance is itself estimated from finite ensembles, and the estimation can make $\chi^{2}$ and each cell's standardized difference larger than this null allows, most along the covariance's smallest-eigenvalue directions, whose estimates are least precise. No correction for this is made, and the verdict stands as graded under this null. \fvval{f1-tune-g18} shares its 2p2h model family with one of the model term's generators. The run-period block was not drawn.
% V-3 (2026-11-16): the paper's approved F2 qualification (DECISION-20261116-joseph-paper-batch.md item 9; paper_body.tex at
% draft/methods-short-20261109 037431430, sec. 6.2's "Their false-flag rates are those of a null ..."), applied to this test.
% absence-ok: "No correction for this is made" scopes the procedure, as the paper's approved sentence does
% absence-ok: scoping statement: effect of omitted covariance terms on flag rate was not measured in OUTCOME-20261108
The result concerns the procedure's covariance at those two tunes, not the reported intervals at the data, where the quoted covariance is made. It is not a coverage rate, says nothing about truths unlike these two or about $C_{\mathrm{PL}}$ or $C_{\mathrm{model}}$, and gives no cause for the joint under-coverage or for any flag. A held-out tune moves the truth, not the sample, so this is not the re-test on held-out simulation of \S\ref{sec:coverage}.
% Note round 1 (2026-11-08, draft for a PRD-referee read and Joseph): OUTCOME-20261108-cov2d-both-tunes-under-cover.md
% @ b551b05f, CITABLE FOR :99-106 (the sentence "At each of two held-out GENIE v3 tunes ... reads UNDER-COVERS", the joint
% count, the flagged cells with pull means and widths, the per-tune chi2/ndf, the false-flag rate), the limits of sec. 4
% :85-95 and the facts sec. 12 requires the text to state (:123-125); NOT CITABLE FOR :110-119 respected (no coverage at
% the data's point, no rate, no cause, nothing on C_PL_R or C_model, "A held-out tune moves the truth, not the sample",
% adds to VL162, no "both COVER"). Every number is \lrval (make_values_landed.py: score_AR23.json fe5263e9,
% score_G18.json 84976694; cell ranges and resolution from the record's lines 52-53 and 91).
% Not drafted here, for a later round (record :127-129, DECISION-20261103-joseph-coverage-clip-rule.md:74-77): the clipping
% rule with C2's and the toys' residuals, outlier cell 216, the weakened cells, the +-1 sigma endpoint draws, the shared
% universes, the cross terms, the single realization, the fixed simulation.

% CV-B04..CV-B09 (sec_systematics.tex:222, sentences 2-4, and the list :225-231 with its comments and :232):
\paragraph{What the covariance does not include.} The covariance gives no calibrated significance for comparisons with generator predictions. We make no claim about coverage, and goodness-of-fit values computed with this covariance are indicative only. The following are not included in the covariance:
\begin{itemize}[nosep]
% (item removed in the D1 follow-on, 2026-11-16: the universe-misses approximation no longer enters the reported covariance; history-cuts/sec_systematics.md S7)
% Search: G2; miss-row weighting; miss-weight-override; universe completeness off/on; G2 repair; stageB/diagnostics P5, stageC/pilot P2-G2, stageC/production, stageF/**; git log --all since 2026-10-15 (only P5 MaRES_0 and P2-G2 Flux:23 at seeds 42/43, both NOT RESOLVED ABOVE TRAINER NOISE: DECISION-20261015-joseph-g2-carry.md:21,27-36; no repair :47) @ 0d919b1a
  \item the dependence on the machine-learning estimator (\S\ref{sec:uq-ml});
  \item the dependence on the regularization, and departures of the detector response from its parametrized family, which we assume absent; we assign no uncertainty for either;
  \item the dependence on the generator shape of the prior, which is reported beside the covariance as $C_{\mathrm{model}}$ (\S\ref{sec:uq-2d-scope}).
\end{itemize}
% J3/J4 provisional (DECISION-20261018, packet W1 register (c): the preparation packet's first two sentences (PR #145), then the director's list of what lies outside every block; the note carries D1's full content, including the reason and "no coverage claim" from D1's must-carry list and the director's draft). Sources: DECISION-20261014-joseph-2d-endpoint-and-open-items.md §2 (construction complete; full scientific closeout NOT ACHIEVED; the three missing inputs; no coverage claim); DECISION-20261017-joseph-2d-adopt-s5-plus-playlist.md §2 (G2 unmeasured; the learner dependence, median 0.98 %; C09 unfolding-model and regularization dependence, C10 response departures: Stage D, unpriced, Stage B REPORT §3 R8). Learner maximum: receipts/2d-lgbm42-20261011/lgbm_central_metrics.out (12.18 %). Review M3: Stage B REPORT §3 (l.313-314) lists C09 (needs independent-generator 2D truths) and C10 ("none outside the parametrized family | stated assumption") separately; only C09 needs independent-generator samples. Review m8 applied. plain-language pass (DECISION-20261018-joseph-plain-language-pass): W1's "Full scientific closeout is not achieved" is carried by "Full validation of the result has not been achieved"; "population or simultaneous interval coverage" by the two coverage cases; "no calibrated covariance-based generator significance" by the next sentence; "not included in any uncertainty block" by "not included in the covariance"; "neither is priced" by "we assign no uncertainty". plain-language pass, review fix m1 (DECISION-20261018-joseph-plain-language-pass): the first two carriers above are superseded. "Full scientific closeout is not achieved" is now carried by "The result's intervals, bias and covariance have not been independently validated", the endpoint ruling's definition (DECISION-20261014 l.49-56: "independent validation that supports the interval, bias and covariance claims"); the coverage cases are "on independent simulated samples" and "simultaneously over all reported bins and projections" (PLAN-20261011 §1: 205 cells, their integral and 30 projections, 236 functionals).

% CV-B10..CV-B14 (sec_systematics.tex:103, with its comments :104-107):
No block contains the dependence on the learner. The central value and the exact-split cross-check differ by a per-bin median of \SI{1.03}{\percent} (maximum \SI{17.6}{\percent}; \S\ref{sec:method-config}). They differ in background treatment as well as in learner. One reseeded exact-split run differs from the unseeded one by a per-bin median of \SI{0.18}{\percent} (maximum \SI{1.6}{\percent}). The spread of the exact-split estimator over seeds is unmeasured (\S\ref{sec:method-config}).
% Review m6, m9. J3/J4 provisional (DECISION-20261018, packet W6: D3 with median and maximum; W9: D11 in
% sec:uq-2d-scope). stageC/production/REPORT.md §3 (C08 4 dof, Gaussian relative SE of variance 71 %; sigma_c08/sigma_ML
% median 0.304; X02 median diag(C_int)/diag(C_stat) 0.0758; interaction/B over c08 diagonal 0.685) and DESIGN.md §8, §12
% (rulings R5, R7). Learner: receipts/2d-lgbm42-20261011/lgbm_central_metrics.out (0.983 %, 12.18 %).
% absence-ok: "No block contains the dependence on the learner" is the source paragraph's own statement (sec_systematics.tex:103 @79a934a41, CV-B10), moved verbatim

% CV-B15..CV-B20 (sec_systematics.tex:275 with its comments :276-282):
\paragraph{Seed-related terms reported, not added.} The variance of the bootstrap\,$\times$\,seed interaction is a per-bin median 0.072 of the statistical block's variance. For the universe\,$\times$\,retraining term, in the six tested universes (of 187) with a non-null shift, the standard deviation of each universe's shift over training seeds, with four degrees of freedom, is a median 1.08--1.36 times the $\sigma$ of the earlier ten-seed block, which is about 3.7 times that of the training-seed block used here. In those universes the universe block, built at one training seed, therefore carries single-seed training noise of about the earlier seed block's size; it is part of the block as built and is not separated from the physical variation. In two of them the noise is comparable to or larger than the shift: in \texttt{2p2h:0} the shift is \SI{0.097}{\percent} against a seed standard deviation of \SI{0.192}{\percent}, and in \texttt{MaCCQE:0} \SI{0.239}{\percent} against \SI{0.210}{\percent} (per-bin medians). We infer that in expectation such noise adds variance to a band rather than removing it, and, from the medians, that these two bands are largely training noise. The universe\,$\times$\,retraining values and the band-centering numbers cannot be recomputed from the saved covariance products; they were therefore not recomputed independently, and their verification is limited to the checksums of their inputs.
% J3/J4 provisional (DECISION-20261018, packet W3: D11; W9: the two universes, note only). stageC/production/REPORT.md §3
% (X02 median diag(C_int)/diag(C_stat) 0.0758; X03 table: 2p2h:0 0.114 % / 0.235 %, MaCCQE:0 0.266 % / 0.220 %, seed SD
% over sigma_ML 1.24-1.33 in the six non-null universes, 4 dof each) and §14 (X03 and centring not recomputable from the
% NPZ); stageC/production/DESIGN.md §8, §12 (rulings R5, R7: reported, not added); DECISION-20261017 §2.
% Review m9: the inference labels follow DRAFT-disclosures D11; F4-VERIFICATION.md compares only the input digests.
% plain-language pass (DECISION-20261018-joseph-plain-language-pass)
% Note register round n3, batch n3-B (2026-10-27): the two inferences share one lead ("We infer"), and the recomputation sentence leads with its reason (R-N3c); the inference labels and every fact are kept.

% CV-B21..CV-B30 (sec_systematics.tex:253 with its comments :254-264):
\paragraph{Inputs carried without confirmation, and playlist pooling.} One systematic category provided in the input products, \texttt{NormDISCC}, is identically zero; we carry it as zero, which we have not been able to confirm is intended. This statement and the proton-on-target statement that follows apply only to the quoted two-dimensional universe block of \S\ref{sec:uq-universes}. The 45-band partition of $C_Z$ (\S\ref{sec:eavailw-construction}), the sixteen-cell covariance (\S\ref{sec:fived16-cov}) and the nuisance draws of the joint tests (\S\ref{sec:joint5d-conditions}) are defined separately. No separate uncertainty is assigned to the counting of protons on target: the NuMI flux prediction carries a \SI{2}{\percent} counting term by convention, so a separate band would risk counting it twice. Two links in that argument, that the flux universes carry the term and that the neutrino--electron constraint absorbs it, remain unverified inferences. The twelve playlists are combined with one global exposure ratio rather than per-period ratios, and the resulting bias is not corrected. At reconstruction level the pooled result is high, by \SI{0.18}{\percent} in total and by at most \SI{0.32}{\percent} in a marginal bin. After unfolding, the pooled result is \SI{0.076}{\percent} above the exposure-weighted combination of the period results on the integral, the same sign as at reconstruction level. It differs from that combination by at most \SI{0.39}{\percent} in a marginal bin (the highest $\pPar$ bin) and by at most \SI{0.33}{\percent} in $\pT$. In single cells the combination differs from the pooled result by \SIrange{-1.83}{3.78}{\percent}, each with $|z|$ at most 0.31.
% J3/J4 provisional (DECISION-20261018, packet W3: D6). docs/publication/DECISION-20261012-joseph-2d-closeout-recipe.md G3;
% DECISION-20261014 item 3 (keep the decision).
% J3/J4 provisional (DECISION-20261018, packet W3: D7). state/2d-closeout/stageB/diagnostics/REPORT.md §6.2 (C07',
% c07prime/DETERMINATION.md; Aliaga et al., arXiv:1607.00704 §IV).
% J3/J4 provisional (DECISION-20261018, packet W3: D9). KNOWN_ISSUES.md J36 (reco level: +0.18 % total, at most +0.32 % in a
% marginal bin; playlist-spread-20261014/REPORT.md §1.2); state/2d-closeout/period-unfold-20261014/REPORT.md (J36
% paragraph) and outputs/analysis.json comparisons.J36. The after-unfolding values print the refined estimator's, regenerated
% by RE-2 part C (stageF/methods/re2/run/partc_r2/partc-r2/values_table.tsv rows pool-integral -0.0764, pool-pz-max
% 0.3948, pool-pt-max 0.3332, pool-cell-min -1.834, pool-cell-max 3.780, pool-cell-zmax 0.308; the "highest p_par bin"
% wording is re2-integration ledger row X1). analysis.json's comparisons.J36 values (-0.0841, -0.5725, +-0.171, -1.793 to
% +3.133, 0.278) are the purity-subtracted estimator's, before RE-2 (n2 review of batch n2-B; comment corrected).

% CV-B31 (sec_systematics.tex:189, sentence 1; the rest stays in the body):
\paragraph{The run-period block $C_{\mathrm{PL}}$ (\S\ref{sec:uq-combine}).} The block's interpretation is limited by the small number of period groups, by possible shrinkage from the regularization, and by the statistical and training noise it retains.

% CV-B32..CV-B34 (sec_systematics.tex:239, sentences 1-3; sentence 1 also stays in the body as the scope label;
% sentences 4-5, the post-measurement identification of the 52 cells, stay in the body):
\paragraph{The model-dependence component $C_{\mathrm{model}}$ (\S\ref{sec:uq-2d-scope}).} $C_{\mathrm{model}}$ measures generator-shape model dependence; it is not a calibrated bound on the bias in each cell. A closure test with NuWro-shaped pseudo-data and the configuration-(a) estimator has a bias above \SI{1}{\percent} in 127 cells. In 67 of them the NuWro shift of the quoted result reaches the bias, and in 72 the quadrature sum of the other two generators' shifts reaches it; these counts are descriptive.
% Source of the 127, 67 and 72 (moved with the text): sec_systematics.tex:234 @79a934a41, the W8 clause.

% Pointer lines (the director's ruling (a), via director-helper-3): caveats kept in the body beside their numbers, listed here.
\paragraph{Comparison with the published result.} The public release cannot establish the origin of the difference between the median of our flux band and that of the published flux block; the FHC-only constraint file is the leading, unconfirmed hypothesis (\S\ref{sec:results-uq}). Beyond the estimator dependence, the source of the shape difference from the published result has not been identified (\S\ref{sec:results-tension}). For binned IBU against the OmniFold result, the origin of the paired-test difference has not been identified, and whether the documented input differences relative to the publication affect the two methods differently was not tested (\S\ref{sec:results-benchmarks}).
% absence-ok: "cannot establish", "has not been identified" and "was not tested" restate sec_results.tex at their homes (sec:results-uq, sec:results-tension, sec:results-benchmarks), which keep them beside their numbers

% Pointer line 4 (the director's ruling (a), via director-helper-3): sec_execsummary.tex:26 / sec:uq-combine, kept there.
\paragraph{Cross-source terms.} The standard covariance sums single-source systematic matrices, each with all other sources nominal, and so omits nonlinear cross-source response terms. No two-dimensional comparison of a joint throw with the block sum is reported (\S\ref{sec:uq-combine}).

% Pointer entry (the director's ruling (a), via director-helper-3): a caveat kept in the body beside its numbers, listed here.


\paragraph{Available energy and momentum transfer.}
% [records reorg] former heading: \subsection{Available energy and momentum transfer (from \S\ref{sec:limitations})}
\label{rlim:3d}
In the comparison with Ascencio \emph{et al.} (\S\ref{sec:4d}), the two analyses define true $\Eavail$ differently, and the aggregate shift over the three lowest bins has not been computed. The definition difference therefore remains an alternative explanation that is not excluded, and at the \SI{10}{\percent} scale of the residual the comparison neither refutes nor validates the result.
% absence-ok: "has not been computed" and "not excluded" restate sec_3d.tex:211, which keeps them beside its numbers
% Entries for the three long caveat statements (the director's ruling, 2026-11-18): each arrives as a move pair from its
% body owner (limitation-resolver: sec_fived16, sec_5d; the Codex writer: sec_joint), whose home keeps a one-line result
% with its scope label and a pointer here. Until each pair lands, its entry is this pointer sentence only.

\paragraph{The sixteen-cell cross section.}
% [records reorg] former heading: \subsection{The sixteen-cell cross section (from \S\ref{sec:limitations})}
\label{rlim:fived16}
% Move pairs LR-20..LR-32 (limitation-resolver): the second coverage test's design and scope, sec_fived16.tex's second-test paragraph (sec:fived16-fd), verbatim and contiguous, except one token: "(above)" becomes (\S\ref{sec:fived16-fd}), where the tunes' overlaps are stated. The home keeps the readings, the counts, "remains untested" and the descriptive sentence.
This is the design and scope of the second coverage test of the sixteen-cell covariance at the two held-out tunes, whose readings \S\ref{sec:fived16-fd} gives. Its pseudo-experiments draw the statistical, flux, interaction and detector uncertainties, with the lateral and normalization ones as surrogates. The graded covariance includes the model term, which is not drawn, since its role is to cover each tune's departure from the simulation. At each tune the covariance is built by the reported procedure using the tune's pseudo-data at their means, with the statistical block bootstrapped from one Poisson realization and the seed block not rebuilt; this fixed covariance stands in for one rebuilt for each pseudo-experiment, which was not computed. The lateral block is carried to each tune in relative form and is not tested, and the tunes' overlaps with the model term's generators (\S\ref{sec:fived16-fd}) apply. The response and run-period terms are left out of the graded covariance. The test has two correlated variants, without and with the GEANT4 neutron and proton bands; the primary variant omits them from both its pseudo-data and its graded covariance. Removing the two bands' block from the covariance graded with them raises the $\chi^2$ of any fixed set of residuals by at most a factor of \fvval{ll5d-mumax-bound-g18} at \fvval{f1-tune-g18} and \fvval{ll5d-mumax-bound-ar23} at \fvval{f1-tune-ar23}. That bound is set by the largest fraction of that covariance, in any direction, that the block carries. The bound applies to the block's place in the covariance only, not to the bands' pseudo-data draws, and what separates the bands' pseudo-experiment and simulation-side blocks remains unclassified. Each tune and each variant is graded separately, per cell and jointly, against thresholds set from the test's null distribution, with a resolution of about \fvval{ll5d-cellres} in coverage per cell and 0.08 jointly. In one variant at one tune, the cell tests raise a false flag with probability at most 0.025, and so does the joint test. These rates are those of a null in which the fixed covariance is exact; the covariance is itself estimated from finite ensembles, which can make $\chi^2$ larger than this null allows, most along its smallest-eigenvalue directions. No correction for this is made. A cell that passes at a tune also passes under the procedure's full covariance there, and the joint reading carries over only if both matrices have 16 degrees of freedom. The model term $V_{\mathrm{prior}}$ of \S\ref{sec:fived16-cov} measures generator-shape model dependence. It is not a calibrated bound. Its coverage at truths held out of its construction has not been validated. The run-period block of the sixteen cells is $\Crun$ (\S\ref{sec:fived16-pl}). In cells \fvval{pl-noise-cells} the block is dominated by the replica noise it contains.
% U11 (currency check of 2fe070fec): the GEANT4 block bound, DECISION-20261116-joseph-paper-batch.md item 3 (A6) wording, from fd/receipts-ll5d/mumax/run/node/mumax.json (422db056; commit e3a47f391), by mumax_values() in make_values_fived16.py.
% absence-ok: "which was not computed", "is not tested" and "No correction for this is made" are scoping statements of the source paragraph, moved verbatim (LR-22, LR-23, LR-31); the declaration fixes them (LL5D-DECLARATION.md :323-324, the home's waiver)
% Move pairs LR-60..LR-62 (limitation-resolver, the structural rewrite): sec_fived16.tex, verbatim; each lead restates its antecedent.
% absence-ok: "is not a calibrated bound" and "has not been validated" are the moved sentences' own statements (LR-60, LR-61; sec_fived16.tex sec:fived16-cov), moved verbatim

\paragraph{The five-dimensional development tests.}
% [records reorg] former heading: \subsection{The five-dimensional development tests (from \S\ref{sec:limitations})}
\label{rlim:5d}
% Move pairs LR-33..LR-47 (limitation-resolver): the coverage test's defective q3 truth (sec_5d.tex sec:coverage-5d), its miscoverage diagnosis (the source block as given, with its comments), and the data-stability consequence (sec:val-departure), verbatim and contiguous. Each lead sentence restates its home's result and supplies the antecedent.
The end-to-end coverage test of \S\ref{sec:coverage-5d} used three truths, and it failed its futility criterion. Of those three, the change of the $q_3$ dependence at fixed $(\Eavail,W)$ did not implement the intended reweighting. Its experiments were generated, and they are not an interpretable test of that deformation. The futility verdict does not depend on them.

\paragraph{Why the intervals miscover under the rate reduction.} The diagnosis below uses the stored outputs and the code.  The rate change itself is recovered: in the one locally stored pseudo-experiment the total changes by \SI{-9.58}{\percent} against \SI{-9.56}{\percent} injected, and over the 400 the total covers at 0.74 (68\%), so the miscoverage is not a normalization failure.  Nor are the intervals too narrow. They use one nominal bootstrap $\sigma$ at every truth, which for 10\% fewer events is slightly too wide apart from the background not being redrawn (inferred spread about $1.15\sigma$ at $W\ge\SI{3}{GeV}$). Setting the spread equal to the half-width leaves the failing quantities' coverage at 0.06--0.22.  The purity subtraction does not apply a nominal signal fraction to the reduced data: it removes the absolute nominal background expectation, which is unchanged, and in the stored pseudo-experiment the subtracted total agrees with the pseudo-signal to 60 of 3.2 million events. The miscoverage is a residual bias that the nominal-truth bias correction does not remove.  The rate change adds a positive bias that grows with $W$: about \SI{0.2}{\percent} for $1.4\le W<\SI{2.2}{GeV}$, \SI{0.4}{\percent} for $2.2\le W<\SI{3}{GeV}$, and 0.7--\SI{1.1}{\percent} for $W\ge\SI{3}{GeV}$. Against intervals of 0.16--0.43\%, this is 1.3--2.6 standard deviations, whatever the sign of the nominal bias $\beta$. Removing that bias restores coverage to 0.62--0.82 at 68\% and 0.92--0.99 at 95\%. These per-quantity biases and spreads are inferred, not measured. The stored output records only hit counts, and each quantity's mean shift and spread come from a Gaussian fit to its 68\% and 95\% coverage, checked against one recorded value (EW17: $+1.04\%$ fitted against $+1.05\%$).  The source of the $W$-graded term is not established. Under an approximate assignment of background fractions to $W$ bins, its size and ordering roughly match the rise of the background fraction by $0.111\,f_{\mathrm{bkg}}$. The term may also arise from the estimator alone. A rerun of five of the same pseudo-experiments did not decide between the two. It was run without background, and with the background scaled by 0.9 so that the background-to-signal ratio is restored. Restoring the ratio raised the residual for $W\ge\SI{3}{GeV}$ rather than removing it, and without background the residual persists. Agreement in the marginal normalizations thus coexists with a substantial residual bias under the tested deformation, and whether that bias has the same origin as the regularization bias of \S\ref{sec:val-departure} is not established. The \estMultiD{} fails the data-stability test (\S\ref{sec:val-departure}). The data bootstrap, which refits the full chain per replica, already includes this rounding-level variability: each replica carries the same spread (median \SI{0.23}{\percent} per joint cell), about \SI{37}{\percent} of the bootstrap variance and below a tenth of any total uncertainty. Central values are therefore reproducible only to within this spread. The conclusion of the comparison with exact binned unfolding (\S\ref{sec:val-departure}) is that the residual is consistent with the regularization of the inverse problem at this binning and prior, and that the tree approximation is not supported as its main cause. This conclusion holds for a binned, signal-only and noise-free comparison at the simulation's detector response, on the development truths, and for these 144 residuals only. It does not address what unbinned or biased estimators could recover, the coverage failure of \S\ref{sec:coverage-5d} or the validity of any interval, and it establishes no impossibility. The precision thresholds that the diagnostic proposes are not acceptance criteria for any result in this note. The maximum real-data prior shifts over five historical vertices are envelopes (\S\ref{sec:gbdt-model-dependence}). These envelopes are sensitivities: they provide no probabilities, covariance components or guaranteed bias bounds, and the truth of the data has not been shown to lie within their span. The studies of \S\ref{sec:gbdt-model-dependence} do not establish a bias--variance tradeoff across estimator settings. Deciding the tradeoff would require sharing pseudo-data and MC splits across a small set of iteration and capacity settings and measuring bias, variance and mean squared error for common functionals. That comparison would not validate intervals, which also requires an explicit treatment of physical nuisances and model dependence. No new ensemble enters these studies. The end-to-end pseudo-experiments with backgrounds return the highest-$W$ column of $(\Eavail,W)$ low at nominal truth (\S\ref{sec:val-negweight}). The bias comes from the purity-based subtraction (\S\ref{app:negweight}), enters no quoted covariance, and has no established size on data. The development assessment of $R$ uses repeated pseudo-experiments (\S\ref{sec:gbdt-model-dependence}). These experiments are not an independent validation: their samples were used during development, and their assessment seeds, although separate, were later also used in the estimator's development.
% [records reorg] former heading: \paragraph{Why the intervals miscover under the rate reduction.}
% Caveat cut list MV13 (DECISION-20261102-joseph-caveat-cut-list.md): the paper's Sec. IV F and its conclusion sentence ("The miscoverage ... is also attributed to a residual bias, but whether that bias has the same origin is not established") moved here; every other fact of IV F was already printed in this section
% Methods-paper rewrite draft 7: docs/orchestration/state/2d-closeout/stageF/methods/M2-RATEFAIL-DIAGNOSIS.md (0 node-h, local receipts and code):
% :8 (total -9.58% recovered against -9.56% injected); :9, :65-69 (spread set equal to the half-width: coverage 0.06-0.22;
% bias removed: 0.62-0.82 at 68%, 0.92-0.99 at 95%); :10, :80-85 (the subtraction removes the nominal absolute
% background expectation, unchanged; the excess is -60 of 3.197M events); :11, :71-76, :100 (W-graded positive residual:
% about +0.2%, +0.4%, +0.7-1.1% for W 1.4-2.2, 2.2-3, >= 3 GeV, 1.3-2.6 sigma against intervals of 0.16-0.43%; added
% whatever the sign of beta); :13, :24 (means and spreads inferred from a Gaussian fit to each 68%/95% coverage pair,
% cross-checked on EW17, +1.04% against +1.05% recorded); :50 (intervals use one nominal bootstrap sigma, slightly too
% wide at 0.9 of the events); :52-63 (failures by W band); :12, :87, :102 (source not established); :15, :104-109
% (confirming test, not run, about 0.1-0.2 node-h).
% 2026-10-22: the test was run (state/ki98-wbias-test/REPORT.md, PR #229; declared and bound before the run), five seeds
% per arm, EW band W >= 3: R_a +2.29 (SE 0.72), R_c +3.40 (0.65), R_b +2.84 (0.46); verdict UNDETERMINED by its rule.
% Move pairs LR-48..LR-50, LR-53 and LR-54..LR-56 (limitation-resolver, the structural rewrite): sec_5d.tex's scope
% sentences, verbatim and contiguous; each lead restates the result that stays in the body and supplies the antecedent.
% absence-ok: "does not address", "establishes no impossibility", "are not acceptance criteria", "provide no probabilities", "would not validate" and "No new ensemble" are the moved sentences' own scoping statements (LR-48..56, sec_5d.tex @ 2fe070fec), moved verbatim
% Move pairs LR-58 and LR-59 (limitation-resolver, the structural rewrite): sec_5d.tex, verbatim; each lead restates its antecedent.
% absence-ok: "enters no quoted covariance", "has no established size on data" and "not an independent validation" are the moved sentences' own statements (LR-58, LR-59), moved verbatim
% Check stage of the journal-voice rewrite (2026-11-21), director-helper-3: the register rows the rewritten Secs. 1, 3 and 6 no longer state.
% [records reorg] % Check stage of the journal-voice rewrite (2026-11-21; CHECK-STAGE.md item 2; DECISION-20261120a item 5), director-helper-3:
% register rows whose limitation the rewritten Secs. 1, 3 and 6 no longer state. Each sentence below is moved verbatim
% from the note as it stood at e89694209 (draft/note-round2-20261108); the source file and line are given with each.
% Under DECISION-20261121b a records-status register row needs only its verbatim text here; the note carries no
% pointer to this section (Joseph, 2026-11-21: note-to-records pointers dropped).
\section{Limitations moved from the note in the journal-voice rewrite}

\subsection{Two-dimensional limitations (from Sec.~3 of the rewritten note)}
From the 2D systematic block's unified-throw test (app\_statmethods.tex:162):
A single-throw superposition probe is insufficient to test nonlinear response.

From the run-period block (sec\_systematics.tex:159):
The run-period block $C_{\mathrm{PL}}$ is an uncertainty assigned for an observed inconsistency between run periods whose cause is unknown. The block does not measure a systematic uncertainty from an identified source.

From the run-period block (sec\_systematics.tex:163):
The block has rank two, and with three period groups its own size is poorly determined.

From the comparison of the flux band with the publication (sec\_results.tex:30):
The public release cannot establish the origin of this difference, for which the FHC-only constraint file remains the leading, unconfirmed hypothesis of this analysis.

From the test inventory's row for the injected $\Eavail$ bump in three dimensions (part5\_validation.tex:56, the result cell):
Marginal residual \SI{1.84}{\percent} of \SI{18.1}{\percent} injected (ratio 0.102); no threshold specified before it ran.

From the table of validation questions for the 2D result (sec\_validation.tex:16, the row for miss completeness):

\begin{tabular}{@{}p{0.2\textwidth}p{0.25\textwidth}p{0.22\textwidth}p{0.25\textwidth}@{}}
Miss completeness & Full truth-event inventory; production truth--reco event matching & Same event keys; truth-only entries retained & Membership and completeness conditions; no recovery or coverage claim \\
\end{tabular}

From the caption of the same table (sec\_validation.tex:7):
Coverage pooled over bins is not a simultaneous probability that every bin is covered.

\subsection{Scope and outlook (from Secs.~1 and 6 of the rewritten note)}
% n01 and n03 (the two introduction sentences below) are retired from the limitations register by
% DECISION-20261121a (Joseph, 2026-11-21); their text stays here.
From the introduction (sec\_intro.tex:11):
Whether it is ready for publication as a measurement has not been established.

From the introduction's scope paragraph (sec\_intro.tex:68):
The results hold for one implementation of unbinned OmniFold, applied to one data set, and for the configurations, observables and deformations tested.

From the outlook (sec\_summary.tex:20):
Further refinements are a fifth external generator, NEUT, and a finer 44-cell comparison with the Ascencio \emph{et al.} measurement (the maximal-common-grid comparison is in \S\ref{sec:4d}).
: its blocks are placed by the routing rules

\paragraph{Two-dimensional limitations.} From the 2D systematic block's unified-throw test (app\_statmethods.tex:162): A single-throw superposition probe is insufficient to test nonlinear response. From the run-period block (sec\_systematics.tex:159): The run-period block $C_{\mathrm{PL}}$ is an uncertainty assigned for an observed inconsistency between run periods whose cause is unknown. The block does not measure a systematic uncertainty from an identified source. From the run-period block (sec\_systematics.tex:163): The block has rank two, and with three period groups its own size is poorly determined. From the comparison of the flux band with the publication (sec\_results.tex:30), where "this difference" is that between the per-bin median of this analysis's flux band, \SI{\fluxBand}{\percent}, and that of the published flux block, \SI{4.01}{\percent}: The public release cannot establish the origin of this difference, for which the FHC-only constraint file remains the leading, unconfirmed hypothesis of this analysis. From the test inventory's row for the injected $\Eavail$ bump in three dimensions (part5\_validation.tex:56, the result cell): Marginal residual \SI{1.84}{\percent} of \SI{18.1}{\percent} injected (ratio 0.102); no threshold specified before it ran. From the table of validation questions for the 2D result (sec\_validation.tex:16, the row for miss completeness):
% [records reorg] former heading: \subsection{Two-dimensional limitations (from Sec.~3 of the rewritten note)}
\begin{tabular}{@{}p{0.2\textwidth}p{0.25\textwidth}p{0.22\textwidth}p{0.25\textwidth}@{}}
Miss completeness & Full truth-event inventory; production truth--reco event matching & Same event keys; truth-only entries retained & Membership and completeness conditions; no recovery or coverage claim \\
\end{tabular}
From the caption of the same table (sec\_validation.tex:7): Coverage pooled over bins is not a simultaneous probability that every bin is covered.

\paragraph{Scope and outlook.} From the introduction's account of the estimators (sec\_intro.tex:59): The \resFiveAxis{} (Table~\ref{tab:config-key}) and its five-axis covariance, and the extended-fiducial value, use the binned purity weight, a per-reco-bin $(D-B)/D$ weight fixed \emph{before} training. From the validity table's row for the exploratory five-axis unfolding (sec\_validity\_table.tex:33): $C_Z$: the earlier note's section on the five-dimensional validation uses it only as a descriptive scale, the \estPurity{}'s total uncertainty, for bias and model-dependence comparisons and for the estimator-seed sensitivity diagnostic (the earlier note's section on the five-axis covariance construction), whose seed-sensitivity measurements use two companion covariances; it also gives the attribution of its per-source decomposition (the earlier note's section on combining the uncertainty blocks). From the introduction (sec\_intro.tex:11), where "it" is the unbinned result in the muon's transverse and longitudinal momenta: Whether it is ready for publication as a measurement has not been established. \emph{Retired (Joseph, 2026-11-21, DECISION-20261121a).} From the introduction's scope paragraph (sec\_intro.tex:68): The results hold for one implementation of unbinned OmniFold, applied to one data set, and for the configurations, observables and deformations tested. \emph{Retired (Joseph, 2026-11-21, DECISION-20261121a).} From the outlook (sec\_summary.tex:20): Further refinements are a fifth external generator, NEUT, and a finer 44-cell comparison with the Ascencio \emph{et al.} measurement (the maximal-common-grid comparison is in \S\ref{sec:4d}).
% [records reorg] former heading: \subsection{Scope and outlook (from Secs.~1 and 6 of the rewritten note)}
% n01 and n03 (the two introduction sentences below) are retired from the limitations register by
% DECISION-20261121a (Joseph, 2026-11-21); their text stays here.
% [records reorg] Joseph, 2026-11-21 (flag c): register row n01 is retired; its text stays, marked
% [records reorg] Joseph, 2026-11-21 (flag c): register row n03 is retired; its text stays, marked
